\snote{settings}

The method depends on a finite-difference step $\delta$, the integrator tolerances, the perturbation size used for admissibility and four thresholds ($\tau_R$, $\tau_\kappa=\tau_{\mathrm{VIF}}$, $\tau_e$). The complete method (selection, admissibility and backward elimination, with the calibration parameters always re-estimated) was re-run changing one setting at a time, with all others at their baseline values ($\delta=1\%$, relative/absolute tolerances $10^{-7}/10^{-10}$, perturbation $\pm5\%$, $\tau_R=0.89$, $\tau_\kappa=\tau_{\mathrm{VIF}}=10$, $\tau_e=0.436$). The 26 problems whose complete method could be re-run fifteen times in reasonable time were used; the other 8 (Alkan, Bachmann, Beer, Chen, Isensee, Lang, Lucarelli, Raimundez) need hours per run. For each setting and problem we record whether the selected subset is identical to the baseline and whether the problem is admissible.

\textit{Being re-computed with the final method for 6 settings; this section is generated automatically when complete.}
